\scanpage{45}
\margin{44}
\uuline{\textcircled{h}} 1) Les faces de $K$ sont indexées par les p.r. p.r. à $\Sigma$ (à l'exclusion des p.r. \deleted{\unsure{exceptionnelles}} supercritiques)

2) Les arêtes de $K$ sont indexées par les \deleted{paires}\inserted{couples} de p.r. \deleted{voisines} $(D,D')$ telles que $D\overset\sigma\longrightarrow D'$ (spécialisations immédiates), \deleted{\unsure{et}} \uline{ou} les flèches des balayages \deleted{entre positions}

\deleted{Si \unsure{les} l'arête $a_\sigma$ est mitoyenne à $F_D$ et $F_{D'}$.}
\deleted{relations voisines d'un}
\deleted{Il est entendu que}

$D\overset\sigma{\text{---}}D'$ entre positions relatives voisines d'un $D_i$ \deleted{$\in\Sigma$} de $\Sigma$ (auquel cas $\unsure{\Sigma|D},\unsure{\Sigma'|D'}$ sont soit stables, soit semi-stables). L'arête $a_\sigma$ est mitoyenne à $F_D$ et $F_{D'}$, et elle est ordinaire ssi $\sigma$ est une spécialisation.

Dans ce cas, le bord privilégié de $a$ au sens de d) est celui qui correspond à la face $F_D$.

3) Les sommets \deleted{\unsure{ordinaires}}\inserted{(ordinaires)} de $K$ sont indexés par les \deleted{triples}\inserted{quadruples} $(D,t',t'',\beta)$, où $D$ est une position relative \inserted{(non supercritique)}, $t,t'$ deux triangles de spécialisation pour $D$ \inserted{(non opposés par un sommet)}, et $\beta$ \inserted{(\unsure{l'un})} un \inserted{des \deleted{quatre}\inserted{deux}} arcs de $\widetilde D$ (un des deux) compris entre les arcs de spécialisation \deleted{\unsure{pour}} correspondant à $t,t'$. (i.e. les \unsure{deux} paires $(D,\beta)$, où $D$ est une pr, et $\beta$ un arc de $\widetilde D$ \deleted{\unsure{limité}} \unsure{déployant|débordant} des deux côtés à un arc de spécialisation ordinaire).

\deleted{Le sommet $s$ associé à cette donnée est exceptionnel ssi $t,t'$ ont \unsure{le} $=$ \unsure{gros un} ``sommet'' de $D$}

(ce qui implique que $D$ est stable).

\deleted{On peut} Les quatre faces incidentes à un sommet $s$ correspondent aux positions relatives $(D,D',D'',D''')$ : $D',D'',\deleted{\unsure{D}}\inserted{D'''}$ se déduisent de $D$ \inserted{par spécialisations}, en balayant $D$ sur $t'$, \deleted{$\unsure{\Sigma}$} sur $t''$, \deleted{et} ou sur les deux à la fois. \unsure{Arêtes} en conséquence

4) Les sommets exceptionnels sont indexés par les repères de $\coprod D_i$, \deleted{\unsure{i.e. des pt}} \deleted{\unsure{les pts}} qui \unsure{revient} au $=$, les \deleted{paires}\inserted{\deleted{couples}} \deleted{\unsure{(s,a) les points de}} \inserted{arêtes orientées de} $\underline\Sigma$, \deleted{\unsure{Au lieu de}} ou les drapeaux $(s,a)$ de type $(0,1)$. \deleted{\unsure{A un tel drapeau}} les \inserted{quatre} faces \deleted{\unsure{issues}}\inserted{issues} de ce sommet sont les quatre positions relatives voisines de $D_i$ \deleted{(celui} qui passe $s$, et qui \deleted{sont} coupent $D_i$ en $s$ resp. en $a$. Les quatre arêtes incidentes en conséquence.

\uline{NB} Ces indications suffisent à construire le \unsure{bâti|tissu} cellulaire $(\mathcal X,K)$, qui est une géométrie d'incidence.

\uline{Cor} L'ens des comp. conn. de $K_0$ est indexé par $I$. La \deleted{$K(i)$} comp. conn. $K(i)$ $(i\in I)$ est canoniquement isomorphe à \deleted{$\unsure{K_0}$} $\widetilde D_i$, l'\unsure{isomo}
